% --- Archon physics formalization source begin ---
% archon:physics
% archon:covers IPhO2026Problems/problem_ipho_2026_t3_a1.lean
% archon:source-report reports/ipho_2026/problem_ipho_2026_t3_a1.source.json
% archon:problem-id ipho_2026_t3
% archon:part-id T3-A1

\chapter{Ipho Physics problem ipho\_2026\_t3\_a1}
\label{ch:IPhO2026Problems_problem_ipho_2026_t3_a1}

\paragraph{Problem source.}
\#\# Physical scenario

A homogeneous isotropic paramagnetic torus has mean radius R, inner radius r
with r << R, volume V, and cross-sectional area A.  An insulated conducting
wire is wound densely around it with N turns and instantaneous current I.
Fields H and B and magnetization M are approximately uniform in the torus.
Use B = mu\_0*H + mu\_0*M, Ampere's law, and the sign convention that work and
heat entering the paramagnetic torus are positive.

\#\# Current subquestion T3-A1

Write the field magnitude H inside the torus in terms of N, I, A, and V.

\paragraph{Shared problem context.}
A homogeneous isotropic paramagnetic torus has mean radius R, inner radius r
with r << R, volume V, and cross-sectional area A.  An insulated conducting
wire is wound densely around it with N turns and instantaneous current I.
Fields H and B and magnetization M are approximately uniform in the torus.
Use B = mu\_0*H + mu\_0*M, Ampere's law, and the sign convention that work and
heat entering the paramagnetic torus are positive.

\paragraph{Current subquestion.}
Write the field magnitude H inside the torus in terms of N, I, A, and V.

\paragraph{Figure/image paths.}
\begin{itemize}
\item ipho\_2026\_source/image/T3\_page-2.png
\item ipho\_2026\_source/image/T3\_page-1.png
\end{itemize}

\paragraph{Formalization target.}
create a compiling Lean file with sorry bodies at `IPhO2026Problems/problem\_ipho\_2026\_t3\_a1.lean`.
The Lean declarations must preserve the physical quantities, dimensions or dimensional roles, figure labels, governing-law hypotheses, and final relation expressed by this problem.
Use Mathlib/Physlib names found through LeanExplore where available. If a domain API is missing, introduce faithful local abstractions rather than scalar placeholder aliases.
This is an answer-blind solve-phase task. Derive a candidate result only from the problem-side evidence above; do not assume a target value, select a tolerance around a desired value, or encode a desired result into a candidate domain.
For numerical questions, define the raw end-to-end quantity and apply an explicit source-derived rounding rule. For identification questions, characterize or prove uniqueness before recording the derived candidate.

\begin{theorem}[Ipho Physics formalization target]
\label{thm:physics:ipho_2026_t3_a1:target}
\leanok
For the densely wound paramagnetic torus in the thin-torus ($r \ll R$)
uniform-field model, the magnitude of the $H$-field inside the torus is
\[
  H \;=\; \frac{N\,I\,A}{V},
\]
in terms of the turn count $N$, the instantaneous current $I$, the
cross-sectional area $A$ and the volume $V$, as requested by T3-A1.  Lean:
\texttt{IPhO2026.T3A1.field\_magnitude} (packaged form
\texttt{IPhO2026.T3A1.field\_magnitude\_eq\_expr}).
\end{theorem}
\begin{proof}
\leanok
Take as the closed curve $C$ of the hint the mean circular path of radius $R$
inside the torus.  In the $r \ll R$ model the azimuthal $H$-field is parallel
to $C$ with constant magnitude $H$, so the circulation evaluates to
$\oint_C \vec H \cdot d\vec\ell = H \cdot (2\pi R)$.  The dense winding threads
the hole of the torus, so the area bounded by $C$ is pierced once by each of
the $N$ turns carrying $I$, giving the linked free current $I_C = N I$.
Ampère's law for the $H$-field therefore reads $H \cdot (2\pi R) = N I$
(\texttt{ampere\_circulation\_solves\_H}).  The torus volume--geometry readout
$V = (2\pi R) \cdot A$ (Pappus; Fig.~3a) multiplied into Ampère's equation
gives $H \cdot V = N I A$ (\texttt{field\_times\_volume}), and division by the
positive volume $V$ yields $H = N I A / V$
(\texttt{field\_magnitude}).  Kernel-checked in Lean with no extra axioms.
\end{proof}
% --- Archon physics formalization source end ---
